% Stage 2 manuscript fragment. Included by master_draft_stage2.tex.
% The adjacent ref/anchor comments are provenance markers for the ARS audit.

\section{Introduction}

Retail recommendation is often presented as a model-selection problem, but an
offline result is first a statement about a defined ranking task. A system must
decide which user history is available, which items are eligible, what counts as
a relevant future event, and how scores are converted into an ordered list. In
large-scale services, candidate generation and ranking are also separate system
roles rather than interchangeable names for one predictor
\cite{covington2016_youtube} %<!--ref:covington2016_youtube--><!--anchor:section:Sections 2 System Overview, 3 Candidate Generation, and 4 Ranking-->
\cite{yi2019_ndr} %<!--ref:yi2019_ndr--><!--anchor:section:Sections 2.3 Two-tower Models, 3 Modeling Framework, and 5 Neural Retrieval System-->.
This distinction matters in retail settings with sparse user--item observations,
changing purchase sequences, item metadata, and products with little or no
training history. The same model label can describe different estimands when
the candidate universe or target event changes.

The present study addresses that problem in a controlled Vietnamese retail
setting. The benchmark is intended to represent the data structures needed for
retail ranking, including user histories, product content, categories, price
signals, and purchase or basket events. Its behavior is controlled and
semi-synthetic rather than a direct sample of Vietnamese shoppers. That choice
allows the protocol to expose specific mechanisms, such as rule-aligned
co-purchase signals and item-side content, while limiting the claims that can be
made about natural behavior. The benchmark declaration specifies 5,000 users,
5,200 items, 823,371 interactions, and a designated item-cold subset; these are
manifest-level design facts and remain subject to the lineage receipt for the
follow-up experiment.

\subsection{Why comparability is a research problem}

Offline recommender comparisons are sensitive to choices that are sometimes
reported as implementation details. Temporal partitioning changes the
information available at prediction time, while target construction determines
which events enter the evaluation population. A recent study of splitting
strategies demonstrates that sequential-recommender outcomes can vary with the
split policy
\cite{gusak2025_time_split} %<!--ref:gusak2025_time_split--><!--anchor:figure:Figure 1; Sections 1–3; global temporal split and results subsections-->
and a systematic replicability analysis of BERT4Rec shows that training and
implementation decisions are part of the reproducibility question
\cite{petrov2022_a_systematic_review_and_replicability_study_of_bert4rec_fo} %<!--ref:petrov2022_a_systematic_review_and_replicability_study_of_bert4rec_fo--><!--anchor:table:Implementation review and RQ2 training-time analysis, including Table 4-->.
Candidate construction, seen-item masking, negative sampling, cutoff, tie
handling, and aggregation consequently define the quantity being estimated.
Metrics with the same name are not automatically comparable when those inputs
differ.

This paper responds by treating the evaluation protocol as a first-class study
artifact. The protocol freezes an immutable temporal split, a complete item
catalog, eligibility rules, relevance semantics, masking, deterministic ranking,
and metric implementation. Each model exports scores to an independent shared
evaluator. The evaluator, rather than the model implementation, applies
masking, ranking, metric calculation, and per-user aggregation. This separation
does not make every method equivalent. It makes the remaining differences
inspectable and prevents a source-native number from silently becoming a result
under a different task contract.

\subsection{Distinct signals and the conditional Hybrid question}

The candidate design combines a deep two-tower scorer with a rule-derived Wide
component. The rationale is a complementarity hypothesis: a rule branch may
retain explicit co-purchase regularities, while a content-aware deep branch may
generalize beyond interactions that have already been observed. This hypothesis
is related to, but not identical with, the Wide \& Deep decomposition of linear
cross-product memorization and embedding-based generalization
\cite{cheng2016_wide_deep} %<!--ref:cheng2016_wide_deep--><!--anchor:section:Sections 2 Recommender System Overview and 3.1-3.3 Wide, Deep, and Joint Training-->.
It is also distinct from item-based collaborative filtering, which aggregates
neighbor evidence from interaction overlap
\cite{sarwar2001_itemcf} %<!--ref:sarwar2001_itemcf--><!--anchor:section:Sections 3.1 Item Similarity Computation and 3.2 Prediction Computation-->,
and from Bayesian Personalized Ranking, which specifies a pairwise objective
for implicit feedback
\cite{rendle2009_bpr} %<!--ref:rendle2009_bpr--><!--anchor:section:Sections 4.1 BPR Optimization Criterion, 4.2 BPR Learning Algorithm, and 4.3.1 Matrix Factorization-->.
The distinction is substantive: a change in encoder, objective, negative
distribution, or system stage can change a result independently of the other
components.

Association rules provide a mechanism vocabulary, not an automatic ranking
claim. Apriori defines support and confidence as thresholds for mining frequent
item relationships
\cite{agrawal1994_apriori} %<!--ref:agrawal1994_apriori--><!--anchor:section:Introduction, formal statement of the association-rule problem and minimum support/confidence thresholds-->,
and multi-item rule systems show how such relationships can be operationalized
for recommendation under a particular dataset and evaluation design
\cite{ghoshal2014_multi_item_rules} %<!--ref:ghoshal2014_multi_item_rules--><!--anchor:section:Abstract; Section 2 problem statement; Section 6 comparison with collaborative filtering and matrix factorization-->.
Prior hybrid and next-basket studies provide useful component-comparison
precedents, including combinations of preference, popularity, transition, and
association signals
\cite{liu2009_hybrid_seq_cf} %<!--ref:liu2009_hybrid_seq_cf--><!--anchor:figure:Replacement work Sections 3 and 3.1; Figures 1–4; Conclusions-->
\cite{peng2022_ham} %<!--ref:peng2022_ham--><!--anchor:table:Modeling sections for long-term preference and association components; Section 6.6 ablation study and Table 13-->
\cite{peng2023_m2} %<!--ref:peng2023_m2--><!--anchor:section:Section 4.1 model components; Section 6.5 ablation study-->,
but those precedents do not determine the contribution of the proposed branch
under the locked v5 cohort.

\subsection{Research question and contributions}

The primary research question is: under a fixed temporal novel-purchase
full-catalog protocol on the controlled Vietnamese retail benchmark v5, does
the proposed decoupled Wide-and-Deep Two-Tower Hybrid improve ranking
effectiveness relative to the strongest faithfully reproduced baseline selected
using validation data? Three subsidiary questions concern item-side cold
performance, the incremental role of the rule branch, and replication on a
compatible public commerce dataset. A fourth concerns efficiency, but it is
opened only after the accuracy artifacts have passed their evidence gates.

The paper makes three contributions that are supported at this stage. First, it
specifies a provenance-aware protocol in which data, split, candidate catalog,
masking, evaluator, seeds, and model artifacts are explicit. Second, it separates
source-native reproduction, public protocol validation, controlled-v5 comparison,
and external validation into distinct evidence namespaces. Third, it reports an
audited public protocol validation lane and preserves its limitations rather
than using it to stand in for the unexecuted Hybrid comparison. The resulting
claim is therefore about reproducible evaluation design and bounded evidence,
not about a settled ordering of recommender architectures. The controlled
Hybrid experiment remains a separately audited follow-up task.

The remainder of the paper first synthesizes prior work by signal and evaluation
contract. It then defines the data and ranking protocol, describes the evidence
lanes, reports the permitted public validation rows, and discusses the limits of
the current evidence.

\section{Related Work}

\subsection{Evaluation and reproducibility}

Recommendation literature often compares model families through headline
metrics, yet the metric is only interpretable together with the prediction
event and candidate set. The split-sensitivity evidence in sequential
recommendation makes temporal construction a substantive design choice
\cite{gusak2025_time_split} %<!--ref:gusak2025_time_split--><!--anchor:figure:Figure 1; Sections 1–3; global temporal split and results subsections-->,
while the BERT4Rec replicability study identifies implementation and training
choices that affect whether a reported result can be reproduced
\cite{petrov2022_a_systematic_review_and_replicability_study_of_bert4rec_fo} %<!--ref:petrov2022_a_systematic_review_and_replicability_study_of_bert4rec_fo--><!--anchor:table:Implementation review and RQ2 training-time analysis, including Table 4-->.
These observations support a narrow methodological conclusion: a comparison
should freeze the target event, temporal boundary, candidate universe, masking,
metric code, aggregation, and training budget before interpreting model
differences.

The present protocol consequently distinguishes source-native evidence from
harmonized evidence. A source-native result retains the source implementation,
data preparation, split, candidate policy, and evaluator. A harmonized result
uses a common data and evaluation contract only after the adapter's deviations
are declared. The distinction protects against a common category error in which
a number reproduced from a paper is treated as if it had been generated under a
new dataset or evaluator. It also makes a failed or incomparable reproduction a
visible outcome rather than a reason to substitute a convenient surrogate.

The positioning is thus deliberately limited. Prior work motivates protocol
control, but it does not supply a ranking for the v5 benchmark or authorize
changing its frozen split after observing outcomes.

\subsection{Collaborative, neural, and sequence recommendation}

Item-based collaborative filtering estimates item similarity from historical
co-interaction and aggregates neighboring items for a user
\cite{sarwar2001_itemcf} %<!--ref:sarwar2001_itemcf--><!--anchor:section:Sections 3.1 Item Similarity Computation and 3.2 Prediction Computation-->.
BPR instead defines a pairwise preference objective for implicit feedback and
can be instantiated with matrix factorization or another differentiable
scorer
\cite{rendle2009_bpr} %<!--ref:rendle2009_bpr--><!--anchor:section:Sections 4.1 BPR Optimization Criterion, 4.2 BPR Learning Algorithm, and 4.3.1 Matrix Factorization-->.
These methods may share user and item identifiers, but they encode different
assumptions about similarity, optimization, and unobserved interactions.

Neural collaborative filtering learns nonlinear interaction functions from user
and item identifiers
\cite{he2017_ncf} %<!--ref:he2017_ncf--><!--anchor:section:Sections 3.1 General Framework, 3.2 GMF, 3.3 MLP, and 3.4 Fusion of GMF and MLP-->.
DeepFM, by contrast, combines low- and high-order interactions over sparse
feature fields for a click-through-rate task
\cite{guo2017_deepfm} %<!--ref:guo2017_deepfm--><!--anchor:section:Section 2 Our Approach, especially 2.1 DeepFM-->.
The shared presence of embeddings does not make their input or target contracts
interchangeable. This distinction informs the adapter registry: a model keeps
its native objective and feature assumptions, while its score output is mapped
to the common evaluator.

Sequential models add another source of non-equivalence. SASRec uses causal
self-attention for next-item prediction
\cite{kang2018_sasrec} %<!--ref:kang2018_sasrec--><!--anchor:section:Section III-C causality; Sections III-E and III-F prediction/objective; Section IV-D evaluation metrics-->,
whereas BERT4Rec uses bidirectional sequence encoding and a masked-item Cloze
objective
\cite{sun2019_bert4rec} %<!--ref:sun2019_bert4rec--><!--anchor:section:Section 3.3 Cloze task; Section 4.2 task settings and evaluation metrics-->.
Their source evaluations use paper-native task and candidate designs. A next-item
objective is therefore not identical to a next-basket set objective, and a
source metric cannot be copied into the v5 table without a protocol bridge.

These families supply comparator and objective precedents, not directly
comparable v5 results.

\subsection{Wide and Deep, Two-Tower, and semantic representations}

Wide \& Deep jointly trains an explicit linear cross-product component for
memorization with an embedding and multilayer component for generalization
\cite{cheng2016_wide_deep} %<!--ref:cheng2016_wide_deep--><!--anchor:section:Sections 2 Recommender System Overview and 3.1-3.3 Wide, Deep, and Joint Training-->.
That decomposition is useful for stating a mechanism hypothesis, but it does not
identify the Wide component with association-rule mining. In the present design,
the rule features are fit from the training transactions and exported as a
separate score. The no-Wide condition is required to assess whether the branch
contributes under the same evaluator.

Two-tower systems are frequently used in large-corpus retrieval. The YouTube
system separates candidate generation and ranking
\cite{covington2016_youtube} %<!--ref:covington2016_youtube--><!--anchor:section:Sections 2 System Overview, 3 Candidate Generation, and 4 Ranking-->,
and neural retrieval work shows that negative-distribution and sampling-bias
choices can be part of the model contract
\cite{yi2019_ndr} %<!--ref:yi2019_ndr--><!--anchor:section:Sections 2.3 Two-tower Models, 3 Modeling Framework, and 5 Neural Retrieval System-->.
DirectAU further illustrates that an objective can target alignment and
uniformity rather than the pairwise objective used by BPR
\cite{wang2022_directau} %<!--ref:wang2022_directau--><!--anchor:section:Sections 2.2 Alignment and Uniformity, 3 Alignment and Uniformity in CF, and 4 Direct Optimization-->.
The protocol must consequently identify whether a tower is a retriever,
pre-ranker, or full-catalog ranker and must record its negative sampler.

Recent tower designs make the same boundary visible from another direction.
ContextGNN introduces pair-aware local graph context while retaining an
exploratory two-tower fallback
\cite{yuan2025_contextgnn} %<!--ref:yuan2025_contextgnn--><!--anchor:section:Abstract and Method sections describing pair-wise familiar-item representations and the two-tower exploratory fallback-->,
and T2Diff adds cross-interaction during training while targeting efficient
matching
\cite{wang2025_t2diff} %<!--ref:wang2025_t2diff--><!--anchor:section:Introduction and Sections 4.2 Diffusion-based Cross Interaction, 4.3 Mixed-Attention Two-Tower, and 4.4 Optimization-->.
These works are useful controls for naming the pair-agnostic boundary. They do
not imply an architecture ordering under the v5 full-catalog contract.

The positioning is therefore specific: this paper evaluates a decoupled
architecture under one locked ranking contract and keeps encoder complexity,
objective, sampling, and system stage as separate explanatory variables.

\subsection{Apriori, association rules, and next-basket recommendation}

Apriori defines a transaction-frequency support threshold and a conditional
confidence threshold for mining association rules
\cite{agrawal1994_apriori} %<!--ref:agrawal1994_apriori--><!--anchor:section:Introduction, formal statement of the association-rule problem and minimum support/confidence thresholds-->.
Multi-item recommendation work demonstrates that such rules can be used to
construct a recommendation list, while its comparisons remain conditional on
the source data and protocol
\cite{ghoshal2014_multi_item_rules} %<!--ref:ghoshal2014_multi_item_rules--><!--anchor:section:Abstract; Section 2 problem statement; Section 6 comparison with collaborative filtering and matrix factorization-->.
The distinction is important for the proposed Wide branch: a mined rule is a
feature or score source, not evidence that the full ranking task has been
solved.

Next-basket research also warns against treating all basket events as one
behavioral regime. A reality-check analysis uses repeat ratios and baseline
behavior to expose how aggregate accuracy can be shaped by repetition
\cite{li2023_nbr_reality} %<!--ref:li2023_nbr_reality--><!--anchor:figure:Figure 1 and Introduction contribution bullets; repeat-ratio and baseline analyses-->,
and exploration-focused work isolates the harder question of recommending novel
items
\cite{li2023_repetition_exploration} %<!--ref:li2023_repetition_exploration--><!--anchor:section:Introduction definitions and research questions; exploration-only analyses-->.
Mask-Swap and its BTBR baseline offer a direct next-novel-basket formulation,
but their reported conditions do not warrant a universal conclusion across
datasets or model families
\cite{li2023_mask_swap} %<!--ref:li2023_mask_swap--><!--anchor:section:Sections 4.1 and 4.2 mask/swap strategy; Section 5 BTBR model; Section 6 experiments-->.
Repeat/explore-aware graph work supplies further component and ablation
precedent under its own data and protocol
\cite{mansouri2026_repeat_explore_lightgcn} %<!--ref:mansouri2026_repeat_explore_lightgcn--><!--anchor:table:Abstract; Table 5 repeat/explore results; Section 6 ablation study; limitations-->.

These studies motivate a train-defined rule-aligned cohort and a no-Wide
ablation. They do not estimate the present mechanism question, and aggregate
basket accuracy alone does not establish a novel-item gain.

\subsection{Graph and contrastive recommendation}

LightGCN simplifies graph convolution by propagating trainable user and item
embeddings over observed interaction edges
\cite{he2020_lightgcn} %<!--ref:he2020_lightgcn--><!--anchor:section:Sections 3.1.1 Light Graph Convolution, 3.1.2 Layer Combination, and 3.3 Model Training-->.
Graph-contrastive controls such as SimGCL and LightGCL add contrastive or
global-relation objectives within their own observed-graph settings
\cite{yu2022_simgcl} %<!--ref:yu2022_simgcl--><!--anchor:section:Sections 2 Empirical Investigation and 3 SimGCL-->
\cite{cai2023_lightgcl} %<!--ref:cai2023_lightgcl--><!--anchor:section:Method section on SVD-guided global collaborative relation modeling and contrastive learning; Experiments-->.
An improvement on a sparse observed graph does not establish a representation
for an item with zero training edges. The v5 protocol therefore reports sparse
and strict zero-edge cohorts separately and does not infer cold-item capability
from a graph score alone.

Graph results motivate comparator coverage and precise cold-item definitions,
but they do not establish v5 performance or zero-edge performance.

\subsection{Content, cold-item, and transfer learning}

Cold-item analysis begins with the training graph. A strict zero-edge item has no
training interaction; an item with one or more training edges may be rare, but it
belongs to a different sparse cohort. DropoutNet formulates cold-start training
through an item-side framework
\cite{volkovs2017_dropoutnet} %<!--ref:volkovs2017_dropoutnet--><!--anchor:section:Section 2, Framework; Section 4.2, Training for Cold Start; Section 5.1, CiteULike-->,
and ALDI defines its own cold-item distillation setting
\cite{huang2023_aldi} %<!--ref:huang2023_aldi--><!--anchor:section:Section 2.1, Problem Formulation; Section 2.2, Existing Solutions-->.
Their thresholds and denominators cannot simply be transferred to v5. A
reproducible cold-item analysis must construct cohorts before scoring and report
warm, sparse, and zero-edge denominators separately.

Text encoders provide a useful item-side representation, but representation
quality is not a ranking outcome. Sentence-BERT supplies a rationale for
semantic sentence embeddings
\cite{reimers2019_sbert} %<!--ref:reimers2019_sbert--><!--anchor:section:Abstract; Section 1, Introduction; Section 3, Training Details-->,
whereas AlphaRec evaluates language representations inside a particular
recommendation and zero-shot protocol
\cite{sheng2025_alpharec} %<!--ref:sheng2025_alpharec--><!--anchor:section:Section 4.1, AlphaRec; Section 5.2, Zero-Shot Recommendation; Section 6, Limitations-->.
Sparse content-based cold-item work likewise defines its own problem and
evaluation setting
\cite{meehan2026_semco} %<!--ref:meehan2026_semco--><!--anchor:section:Section 1, Introduction; Section 3.1, Problem Definition; Sections 3.2-3.3-->.
The proposed pipeline therefore requires a recommender objective and a shared
task-specific evaluator in addition to an embedding recipe.

Transferable item representations retain target-domain assumptions. UniSRec
includes downstream adaptation
\cite{hou2022_unisrec} %<!--ref:hou2022_unisrec--><!--anchor:figure:Figure 1; Sections 2.2-2.4, especially 2.4 Adaptation to Downstream Recommendation Tasks-->,
VQ-Rec evaluates vector-quantized transfer under its experimental setup
\cite{hou2023_vqrec} %<!--ref:hou2023_vqrec--><!--anchor:section:Sections 2.1-2.3; Section 3.1 Experimental Setup-->,
and UTGRec includes downstream fine-tuning
\cite{zheng2026_utgrec} %<!--ref:zheng2026_utgrec--><!--anchor:section:Section 2.4.2, Downstream Fine-tuning; Sections 3.1.1 and 3.1.3-3.1.4; Section 3.3, Ablation Study-->.
An architecture inspired by these works is not a replication unless the encoder,
data lineage, adaptation procedure, split, candidates, metrics, and statistics
are matched or every deviation is declared. Shared platform or benchmark
lineage also means that pretraining and downstream results cannot be counted as
independent replications merely because their domain labels differ
\cite{meehan2025_cold_popbias} %<!--ref:meehan2025_cold_popbias--><!--anchor:section:Section 1, Introduction; Sections 3.1-3.3-->
\cite{meehan2026_semco} %<!--ref:meehan2026_semco--><!--anchor:section:Section 1, Introduction; Section 3.1, Problem Definition; Sections 3.2-3.3-->.
Content aligned to collaborative teachers can also inherit popularity structure,
so content availability alone does not guarantee an unbiased or effective
cold-item ranking
\cite{meehan2025_cold_popbias} %<!--ref:meehan2025_cold_popbias--><!--anchor:section:Section 1, Introduction; Sections 3.1-3.3-->.

The positioning is conditional on full contract alignment. External and
cold-item conclusions require matched encoder, data, adaptation, split,
candidate set, metric, and statistical procedure; a source architecture name is
not sufficient evidence.

